\section{Beyond Locks}

\subsection{Deadlocks}

\begin{remark}
\textbf{Deadlock detection} in systems is implemented by finding cycles in the dependency graph.
\end{remark}

\begin{remark}
\textbf{Deadlock avoidance} techniques to ensure a cycle can never arise:
\begin{itemize}
    \item \textbf{Imposing an ordering} on the resources and requiring processes to acquire resources in that order
    \item \textbf{Two-phase locking protocol} with retry and release on failure
\end{itemize}
\end{remark}

\begin{codeexample}
//Programming trick to create a global ordering
class Resource {
    private static final AtomicLong counter = new AtomicLong();
    private final long index = counter.incrementAndGet();
    ...
}
\end{codeexample}

\subsection{Semaphores}

\begin{definition}
A \textbf{semaphore} is a synchronization primitive that acts as a thread-safe counter to control access to a shared resource by multiple processes.
\end{definition}

\begin{codeexample}
public class Semaphore {
    private int n, S;
    acquire() {
        wait until S > 0;
        S--;
    }
    release() {
        S++;
    }
}
\end{codeexample}

\begin{codeexample}
Semaphore semaphore = new Semaphore(n);
semaphore.acquire();
    //no more than n threads can be here at once
semaphore.release();
\end{codeexample}

\begin{remark}
Semaphores are a generalization of locks and can be used to implement locks and other synchronization primitives.
\end{remark}

\begin{definition}
A \textbf{binary semaphore} is a semaphore that can only have values 0 and 1, this can be used to implement locks.
\end{definition}

\begin{definition}
A \textbf{rendezvous} is a location in the code where threads have to wait until all threads have reached that location before they can continue.
\end{definition}

\begin{remark}
For two threads, a rendezvous can be implemented with two binary semaphores.
\end{remark}

\begin{codeexample}
void thread1() {
    A.acquire();
    B.release();
    B.acquire();
    A.release();
}
void thread2() {
    B.acquire(); 
    A.release();
    A.acquire();
    B.release();
}
\end{codeexample}

\begin{important}
Using Semaphores is often problematic as the programmer has to have deep knowledge of the algorithm to correctly implement the synchronization.
\end{important}

\subsection{Barriers}

\begin{definition}
A \textbf{barrier} is a synchronization primitive that coordinates threads to wait until all have reached the barrier, then releases them all at once.
\end{definition}

\begin{definition}
A \textbf{two-phase barrier} is implemented with two semaphores and a counter.
\end{definition}

\begin{codeexample}
mutex=1, barrier1=0, barrier2=1, count=0;

aquire(mutex);
    count++; 
    if (count == n) {
        accuire(barrier2);
        release(barrier1);
    }
release(mutex);

acuire(barrier1); release(barrier1);

aquire(mutex);
    count--;
    if (count == 0) {
        accuire(barrier1);
        release(barrier2);
    }
release(mutex);
\end{codeexample}

\subsection{Producer Consumer Pattern}

\begin{definition}
The \textbf{producer consumer pattern} is a synchronization pattern where producers generate data and consumers process it, coordinated through a shared buffer.
\end{definition}

\begin{remark}
We can use producer-consumer patterns to build data flow parallel programs (pipelines):
$T_0 \rightarrow T_1 \rightarrow T_2 \dots T_n$
where $T_0$ is a producer, $T_n$ is a consumer, and intermediate nodes are both.
\end{remark}

\begin{remark}
A \textbf{Producer-Consumer Queue} with buffer improves throughput. We stall $T_i$ iff the queue $Q_i$ is full. Implemented using Semaphores that signal producers (space available) and consumers (data available).
\end{remark}

\begin{codeexample}
void enqueue(long x) {
    try {
        nonFull.acquire();
        manipulation.acquire();
        buffer[in] = x;
        in = (in + 1) % n;
    } catch (Exception e) {}
    finally {
        manipulation.release();
        nonEmpty.release();
    }
}
\end{codeexample}

\begin{codeexample}
long dequeue() {
    long x = 0;
    try {
        nonEmpty.acquire();
        manipulation.acquire();
        x = buffer[out];
        out = (out + 1) % n;
    } catch (Exception e) {}
    finally {
        manipulation.release();
        nonFull.release();
    }
    return x;
}
\end{codeexample}

\subsection{Monitors}

\begin{definition}
A \textbf{monitor} is a synchronization primitive that allows threads to wait on condition variables to be notified about changes in state.
\end{definition}

\begin{definition}
A \textbf{condition variable} is a condition that threads can wait on. Other threads can signal it to notify waiting threads.
\end{definition}

\begin{codeexample}
Monitor buffer = new BufferMonitor(10);
Condition nonFull = buffer.newCondition();
Condition nonEmpty = buffer.newCondition();

void producer() {
    buffer.lock();
    try {
        while (buffer.isFull()) { nonFull.await(); }
        buffer.enqueue(i);
        nonEmpty.signal();
    } finally { buffer.unlock(); }
}

void consumer() {
    buffer.lock();
    try {
        while (buffer.isEmpty()) { nonEmpty.await(); }
        buffer.dequeue();
        nonFull.signal();
    } finally { buffer.unlock(); }
}
\end{codeexample}
